\subsection{Псевдокод алгоритма}

Команда \verb|\algo| предназначена для оформления вычислительных процедур.
Внутри неё можно использовать окружение \verb|algorithmic|, а в необязательных аргументах указать вычислительную сложность и замечание.

\algo[optional=note]{algo:сертификация_агента}
{Сертификация агента-кандидата}{
    \begin{algorithmic}[1]
        \Require показатели $p, c, h \in \Nz$
        \Ensure решение о сертификации
        \State $r \gets 2p + c - 3h$
        \If{$h = 0 \land r \ge 10$}
            \State \Return <<допустить>>
        \Else
            \State \Return <<отправить на коррекцию>>
        \EndIf
    \end{algorithmic}
}[
    В модели с арифметическими операциями постоянной стоимости алгоритм выполняется за $\bo{1}$ операций.
][
    Настоящая процедура допуска системы искусственного интеллекта не должна состоять из одной проверки условного оператора.
    Даже если этот оператор записан достаточно убедительным псевдокодом.
]

Алгоритм вычисляет рейтинг из \blink{определения рейтинга}{defi:рейтинг_агента} и проверяет условия из \blink{определения сертифицированного агента}{defi:сертифицированный_агент}.


\subsection{Задачи с отображаемыми решениями}

Глобальная команда \verb|\showtaskdetailstrue| включает отображение решений и ответов.

\showtaskdetailstrue

\task{task:чётность_рейтинга_агента}{
    Пусть агент-кандидат имеет показатели $A = \tup{p, c, h}$, причём числа $c$ и $h$ имеют одинаковую чётность.
    Докажите, что рейтинг $R(A)$ является чётным числом.
}[
    По \blink{определению рейтинга}{defi:рейтинг_агента} имеем $R(A) = 2p + c - 3h = 2(p - h) + (c - h)$.
    Число $2(p - h)$ чётно, а число $c - h$ чётно вследствие одинаковой чётности $c$ и $h$.
    Следовательно, рейтинг $R(A)$ является чётным числом.
]

\task{task:число_раундов_коррекции}{
    Агент имеет начальные показатели $A_0 = \tup{1, 2, 2}$.
    Определить минимальное число описанных выше раундов коррекции, после которого агент может быть сертифицирован.
}[
    До устранения всех галлюцинаций агент не может быть сертифицирован, поэтому требуется не менее двух раундов.
    После двух раундов получаем $A_2 = \tup{5, 4, 0}$ и $R(A_2) = 14$.
    Следовательно, двух раундов достаточно.
][
    2 раунда.
]

Если решение отсутствует, но требуется вывести ответ, перед ним указывается пустой необязательный аргумент.

\task{task:минимальное_число_задач}{
    Агент имеет показатели $A = \tup{p, 4, 0}$.
    Найти наименьшее значение $p$, при котором агент сертифицирован.
}[][
    3.
]


\subsection{Задачи со скрытыми решениями}

Команда \verb|\showtaskdetailsfalse| скрывает решения и ответы, но сохраняет их в исходном файле.
Такой режим подходит для версии документа, предназначенной студентам.

\showtaskdetailsfalse

\task{task:компенсация_галлюцинации}{
    Агент имеет показатели $A = \tup{4, 5, 1}$.
    Сколько дополнительных правильно решённых задач необходимо добавить без изменения остальных показателей, чтобы рейтинг стал не меньше $10$?
    Будет ли после этого агент сертифицирован?
}[
    Начальный рейтинг равен $R(A) = 2 \mul 4 + 5 - 3 = 10$.
    Поэтому для достижения рейтинга $10$ дополнительные задачи не нужны.
    Однако агент не будет сертифицирован, поскольку $h = 1$.
][
    Дополнительные задачи не требуются, но агент не будет сертифицирован.
]

В текущей версии документа решение и ответ последней задачи не отображаются.
Для подготовки преподавательской версии достаточно заменить команду \verb|\showtaskdetailsfalse| командой \verb|\showtaskdetailstrue| или централизованно включить соответствующий флаг перед началом документа.


\subsection{Программный код и вывод}

Команда \verb|\code| подключает исходный файл на языке Python.
Необязательный четвёртый аргумент задаёт файл с выводом программы.
Вызовы функции \mintinline{python}{print} можно сопоставить со строками вывода с помощью маркеров \verb|\cmo|.

\code
    {code:сертификация_агентов}
    {Сертификация агентов на языке Python}
    {code/l-02/code1.py}
    [code/l-02/code1.txt]

\blink{Листинг программы сертификации}{code:сертификация_агентов} является реализацией \blink{алгоритма сертификации}{algo:сертификация_агента}.
Бейджи справа от вызовов \mintinline{python}{print} указывают номера соответствующих строк в блоке вывода.